\section{Honest AI Claims}\label{sec:honest_ai_claims}

This section is a deliberate, dedicated audit of what the v0.4 evidence does and does not let us say about ``creativity'' in PCE. Three contemporary failure modes in AI papers about creative generation motivate the section: (a) overclaiming a quantitative win on a small benchmark slice that is then extrapolated to ``LLMs are creative''; (b) underclaiming by retreating to ``the architecture is interesting'' once the headline contrast fails to land; (c) silently swapping the operationalisation of ``creativity'' between the introduction and the discussion. We try to do none of those.

\subsection{What ``creative'' means here, operationally}

Throughout this paper, ``creative output'' is operationalised as the per-item composite score under the per-domain rubric documented in \S\ref{sec:benchmarks} and Appendix~\ref{app:scoring}. The composite is a weighted sum of axis scorers (creativity, lexical diversity, idiosyncrasy, emotional resonance, literary devices, imagery for poetry-gen; aspect-multiplicity for poetry-interp; novelty + framing-coverage + depth for sci-creativity; CreativityPrism Q$\cdot$N$\cdot$D for AUT). The composite is a learned proxy for human inter-rater quality, not a definition of creativity in any substantive sense. When we say ``the cascade produced a better surface'' we mean ``the composite scored the surface higher''; when we say ``the judge disagrees with the scorer'' we mean ``the Sonnet-4 judge's per-axis scoring on the same surfaces does not correlate with the Haiku composite''. The substantive question of whether a higher composite is the right notion of creativity is open and is not what this paper claims to settle.

\subsection{What v0.4 does \emph{not} claim}

v0.4 does not claim:

\begin{enumerate}[leftmargin=2em,itemsep=2pt,topsep=2pt]
\item that PCE makes Claude Haiku more creative in any human-judged sense (we have not yet run a human-evaluator study at scale; the v0.5 ladder includes one);
\item that the cascade beats bare Haiku at the headline cascade-vs-bare contrast (H1.v4--H4.v4 are inconclusive at the available $n$);
\item that the \iast{vimar\'sa} event-gate is a working insight detector (H8b shows it is mis-calibrated);
\item that the small Haiku composite-score head agrees with the larger Sonnet-4 judge (H9 shows it does not, on the cascade items);
\item that the architecture has been ablated component-by-component (the Hopfield-warm-start signal alone, the BMR aspect-conditioning alone, etc., are deferred to v0.5);
\item that the Sanskrit chandas showcase (\S\ref{sec:showcase}) is a cascade-generated artefact (it is maintainer-curated reference text validated by the chandas tool).
\end{enumerate}

\subsection{What v0.4 \emph{does} claim}

v0.4 does claim, and the evidence supports:

\begin{enumerate}[leftmargin=2em,itemsep=2pt,topsep=2pt]
\item the cascade's second pass reliably produces a higher-composite surface than the first ($g=+0.65$, $p\sim 10^{-4}$, $n=27$; H8a);
\item among the four commit policies tested, \texttt{always\_revise} dominates \texttt{learned\_gate} dominates \texttt{event\_gated} dominates \texttt{always\_draft}, with statistically significant pairwise gaps at the extremes (H8c);
\item a learned commit gate trained on the H8a paired-revision data improves over a hand-coded \iast{vimar\'sa} disjunctive event signature (F1 = 0.65 vs.\ 0.52; H8b);
\item the Sonnet-4 judge and the Haiku composite scorer do not agree on the sign of the cascade-vs-bare delta on the v0.4 cascade items, at a rate (56.5\% sign-agreement) consistent with random ($\rho=0.0$, $p=1.0$; H9);
\item the v0.4 cost ledger is internally consistent, externally reproducible (\$12.73 over 1{,}277 Bedrock calls), and traceable to the per-call accounting in \texttt{audit/v0.4/cost\_ledger\_merged.json};
\item the v0.4 plugin and the standalone CLI are byte-for-byte the same substrate (\S\ref{sec:substrate}), so a third-party reproduction does not depend on an IDE plugin runtime.
\end{enumerate}

\subsection{Why we keep the Pratyabhij\~n\=a framing}

The Pratyabhij\~n\=a tradition is a real philosophical lineage with millennia of internal debate over the operators we appropriate. It would be possible (and, in a paper that wanted to maximise reception in mainstream ML venues, advisable) to strip the Sanskrit terminology and present PCE as ``a draft-shadow-revision-commit-multiplex cascade with a recursive critique step gated on a free-energy proxy''. We deliberately do not do this. The Pratyabhij\~n\=a vocabulary is what made the gating-and-pruning question visible to us in the first place; the contemporary self-refinement literature is converging on the same questions but did not, when this work began, have a precise enough vocabulary to ask them in the form ``\emph{which \iast{\'sakti} should fire next?}''. We owe the philosophical lineage the citation \citep{Lawrence1996TantricArgumentPEW,abhinavagupta-isvarapratyabhijna,pratyabhijnahrdayam-singh}; we do not owe the lineage the metaphysics, and we explicitly do not import it.
